% ==========================================================================
% AAAI-submission draft (self-contained; compiles with pdflatex).
% For camera-ready: replace the preamble block below with \documentclass[letterpaper]{article}\usepackage{aaai2026}
% and the standard AAAI boilerplate. Content/section structure already matches AAAI.
% ==========================================================================
\documentclass[twocolumn,10pt]{article}
\usepackage[letterpaper,margin=0.75in]{geometry}
\usepackage{times}
\usepackage{amsmath,amssymb}
\usepackage{booktabs}
\usepackage{graphicx}
\usepackage[hidelinks]{hyperref}
\usepackage{xcolor}
\usepackage{caption}
\captionsetup{font=small}
\setlength{\parskip}{0pt}
\frenchspacing

\title{\bf When Does a Learned World Model Help?\\
A Value-Sufficiency Diagnostic that Predicts\\ World-Model Dependence Across Architectures}
\author{Anonymous submission}
\date{}

\begin{document}
\maketitle

\begin{abstract}
Whether a learned world model (WM) helps a model-based agent is usually
debated at the level of \emph{architectures}: latent-consistency (TD-MPC2)
versus reconstruction RSSM (Dreamer) versus reward-only value learning.
We argue the more predictive question is at the level of \emph{tasks}.
We introduce the \textbf{Value-sufficiency BottleNeck (VBN)}, a cheap,
checkpoint-time probe that measures how much of an agent's latent the value
function actually needs, and show that this single scalar \emph{predicts}
whether ablating the world model's forward-dynamics learning helps or hurts
return. Across three DeepMind Control tasks and two independent WM families
(Dreamer, TD-MPC2) the same task ordering holds: on \textsc{acrobot-swingup}
the WM is essential (removing it collapses return by $6.8\times$; $+311$ to
$+1458\%$ per seed), on \textsc{cheetah-run} it helps ($+11\%$), and on
\textsc{walker-run} it is redundant (stripped $\ge$ vanilla). This ordering
matches each task's VBN fingerprint. We also report a clean \emph{negative}:
a plan-time ``gate'' that down-weights the WM's imagined rollout---the
prescriptive fix the diagnostic seems to suggest---does \emph{not} robustly
beat trusting the full model ($n{=}5$; all gate settings inside a 19-point
band). Our contribution is therefore a \emph{diagnostic}, not a repaired
algorithm: a task property, groundable in the value-equivalence principle
and rate--distortion, that tells you \emph{before} you spend the compute
whether a learned world model has room to help.
\end{abstract}

\section{Introduction}
Model-based reinforcement learning couples a learned \emph{world model}
(WM)---a latent dynamics model trained to predict transitions---with a value
function and a planner. A large literature asks whether this pays off, and
answers by comparing architectures: generative reconstruction models
\cite{hafner2023dreamerv3}, latent-consistency models
\cite{hansen2024tdmpc2}, and value-equivalent models
\cite{schrittwieser2020muzero,grimm2020value}. The empirical picture is
mixed, and the mixedness is usually attributed to modeling choices.

We take a different stance. We hold the architecture fixed and ask which
\emph{tasks} make the learned dynamics load-bearing. Concretely, we ablate
the WM by disabling its forward-dynamics learning objective---in TD-MPC2 by
setting the latent-consistency coefficient to zero, in Dreamer by zeroing the
dynamics loss---leaving the value function, encoder and planner otherwise
intact. We call the change in return the task's \emph{WM-dependence}.

Our central finding is that WM-dependence is predictable from a task property
we can measure cheaply and \emph{a priori} of the ablation. The property is
how compressible the value-relevant information is: if a tiny bottleneck on
the value head already recovers most of the return, the value function does
not need the full latent, and imposing accurate forward dynamics on that
latent is redundant. If the value head needs most of the latent, the forward
model's job of shaping and propagating that information becomes load-bearing.
We operationalize this with the Value-sufficiency BottleNeck (VBN).

\paragraph{Contributions.}
\begin{enumerate}\itemsep2pt
\item \textbf{VBN}, a checkpoint-time probe of value-relevant
compressibility, and its three-shape task fingerprint (Sec.~\ref{sec:vbn}).
\item A \textbf{cross-model, cross-task} result: the same three tasks flip
between WM-essential, WM-helpful, and WM-redundant in \emph{both} Dreamer and
TD-MPC2, and the ordering matches the VBN fingerprint (Sec.~\ref{sec:cross}).
\item An honest \textbf{negative}: the ``gated world model'' that the
diagnostic seems to prescribe does not robustly beat the full model
(Sec.~\ref{sec:gate}).
\item A \textbf{mechanism}: under planner-driven data collection the WM can
actively \emph{hurt} by inflating the planner's own target variance
(Sec.~\ref{sec:mechanism}).
\item A \textbf{theoretical grounding} of VBN in the value-equivalence
principle and rate--distortion, and a roadmap toward a training-free
analytical criterion (Sec.~\ref{sec:theory}).
\end{enumerate}

\section{The Value-sufficiency BottleNeck}
\label{sec:vbn}
Let $z=\phi(o)\in\mathbb{R}^{d}$ be the agent's latent and $V_\theta(z)$ its
value head. VBN inserts a linear bottleneck of width $D\le d$ into the value
head: $\hat V(z)=g_\psi(W_2\,\sigma(W_1 z))$ with $W_1\in\mathbb{R}^{D\times d}$,
and re-trains only the small heads while freezing $\phi$. Sweeping
$D\in\{16,32,64,128\}$ and recording the resulting control return traces a
discrete \emph{value rate--distortion} curve: how much return survives when
the value function is forced through $D$ dimensions.

Empirically the curve has one of three shapes, which we treat as a task
fingerprint:
\begin{itemize}\itemsep2pt
\item \textbf{monotone} (return keeps rising with $D$): value needs a large
fraction of the latent---low compressibility.
\item \textbf{flat-high} (already saturated at small $D$): a tiny bottleneck
suffices---high compressibility, value is nearly sufficient in few dimensions.
\item \textbf{ramp-to-$D^\star$} (rises then plateaus at an intermediate
$D^\star$): value needs a specific, sizeable subspace.
\end{itemize}
In our four-task grid ($n{=}3$ seeds each), \textsc{cheetah-run} is monotone,
\textsc{walker-run} is flat-high, \textsc{acrobot-swingup} ramps to
$D^\star{=}64$, and \textsc{hopper-hop} is flat-high/noisy. VBN costs one
short probe-training run per width on an existing checkpoint; it never
touches the base policy.

The claim we test is: \emph{low value-compressibility (monotone / ramp)
predicts that the learned WM is load-bearing; high compressibility
(flat-high) predicts the WM is redundant.}

\section{Cross-Model, Cross-Task Result}
\label{sec:cross}
We ablate the WM in each framework and measure the change in return per task.
For Dreamer we use the unmodified official DreamerV3 and report the
last-30-episode median return at $\sim$1.1M environment steps. For TD-MPC2 we
use our JAX reimplementation, audited against the official release: hopper-hop
is at parity (within the official run-to-run seed spread), cheetah $-5\%$,
walker $-17\%$, acrobot $-23\%$, with known causes (data-collection mode,
physics backend) under verification. All TD-MPC2 results here are
\emph{within-implementation relative} comparisons (vanilla vs.\ ablated under
identical code), so these absolute gaps do not affect the claims; the
architecture-general headline rests on unmodified Dreamer. The strip ablation
removes only the forward-dynamics learning signal.

\begin{table}[t]
\centering\small
\begin{tabular}{lccl}
\toprule
Task & VBN & Strip-WM effect (Dreamer) & Verdict \\
\midrule
\textsc{acrobot} & ramp$\to$64 & van $412.5 \gg$ strip $60.8$ & WM \emph{essential} \\
\textsc{cheetah} & monotone    & van $710 >$ strip $637$      & WM \emph{helps} \\
\textsc{walker}  & flat-high   & strip $776 \ge$ van $722$    & WM \emph{redundant} \\
\bottomrule
\end{tabular}
\caption{The cheap VBN fingerprint predicts the WM-dependence ordering,
$n{=}3$ seeds each. Acrobot: $6.8\times$ gap ($+311$ to $+1458\%$ per seed).
Cheetah: $+9.7$--$11.4\%$ across seeds. Walker: null (stripped mildly higher).
The same ordering holds in TD-MPC2 (latent-consistency): cheetah
WM-load-bearing, walker redundant.}
\label{tab:cross}
\end{table}

Table~\ref{tab:cross} is the core result. Three points deserve emphasis.
\textbf{(i)} The ordering is \emph{monotone in the VBN fingerprint}: the more
of the latent the value head needs, the more removing the WM hurts.
\textbf{(ii)} The ordering is \emph{architecture-invariant}: it reproduces
across a reconstruction RSSM (Dreamer) and a latent-consistency model
(TD-MPC2), two families that share almost no inductive bias except that they
both learn forward dynamics. \textbf{(iii)} The high-dependence anchor is
\emph{dramatic but honest}: on acrobot, vanilla return is remarkably stable
across seeds ($408,420,409$) while the \emph{stripped} runs are high-variance
($26,56,100$)---sometimes near-total failure, sometimes a partial recovery.
The qualitative claim (the WM is essential on acrobot) is unshakable; the
magnitude of the collapse varies by seed. We report this rather than hiding
the variance.

\section{A Clean Negative: The Gate Does Not Help}
\label{sec:gate}
The diagnostic appears prescriptive: if we can measure, online, that a task
(or state) is WM-redundant, we should be able to \emph{down-weight} the WM
there and win. We implemented exactly this as a plan-time \emph{gate}
$g\in[0,1]$ that scales the planner's weight on the WM's imagined rollout,
$\text{score}=g\!\sum_t \gamma^t r^{\text{wm}}_t+\gamma^H V$, so that $g{=}1$
is the full model and $g{\to}0$ trusts the imagined rollout less.

On \textsc{cheetah-run} under planner-collection, an early single-seed run
suggested a modest win for $g{=}0.25$ ($+5.6\%$). It did not survive
replication. At $n{=}4$--$5$ seeds the per-gate means collapse into a
19-point band:
\[
\underbrace{701.9}_{g=0.0}\;
\underbrace{688.4}_{g=0.5}\;
\underbrace{686.5}_{g=1.0}\;
\underbrace{682.9}_{g=0.25},
\]
while the within-gate seed spread is $\sim$80--100 points (e.g.\ $g{=}0.25$
spans $637$--$715$). No gate value robustly beats the full model; the apparent
$n{=}1$ advantage was seed noise. \textbf{The tunable gate is a confirmed
negative.}

\paragraph{Why the null is structural, not mysterious.}
Post-hoc analysis shows the gate as instantiated could hardly have mattered.
With horizon $H{=}3$ and $\gamma{=}0.99$, the score decomposes as
$g\sum_{t<3}\gamma^t r_t + \gamma^3 V \approx 2.97\,\bar r\;(\text{gated}) +
97\,\bar r\;(\text{terminal})$: the gated term is $\sim$3\% of the trajectory
score, so sweeping $g\in[0,1]$ perturbs the planner objective by at most a few
percent---below seed noise by construction. Moreover $g{=}0$ is \emph{not}
WM-free planning: the terminal value $V(z_H)$ is still evaluated on a latent
rolled through the learned dynamics, and the value function was trained with
the WM regardless; the gate removes imagined \emph{reward}, not imagined
\emph{state}. Correct instantiations of the same idea---horizon gating
($H\!\to\!0$ yields value-greedy planning), per-step uncertainty weighting in
the spirit of STEVE \cite{buckman2018steve}, or training-time gating---remain
untested future work. We keep the negative because it clarifies the
contribution: this paper offers a \emph{diagnostic}, not a repaired algorithm,
and knowing \emph{whether} a WM will help is separable from---and, on this
evidence, easier than---\emph{fixing} one that does not.

\section{Mechanism: How the WM Can Hurt}
\label{sec:mechanism}
Redundancy alone would predict a null, not a loss. Yet on \textsc{cheetah}
under \emph{planner-driven} collection (data gathered by the MPPI planner
rather than the policy), \emph{removing} the WM \emph{helps} by $+45\%$
($n{=}9$)---an inversion. We trace it to variance: planner-collection inflates
the variance of the planner's own value target by $\approx 3\times$ (the
``H-variance'' effect). The imagined rollout, when the value landscape is
already nearly sufficient in few dimensions, adds a high-variance term to the
plan-time objective that the planner then chases, destabilizing learning.
This is the concrete failure mode the diagnostic flags: the WM does not merely
fail to help; it can inject target variance where the value function did not
need it.

\section{Grounding: Compressibility, Value-Equivalence, and an Analytical Criterion}
\label{sec:theory}

\paragraph{Is ``compressibility'' well-defined?}
Yes, information-theoretically---and VBN is a cheap estimator of it, not the
exact object. The quantity we care about is the minimal description of state
that suffices to represent value. Three standard formalizations coincide here:
(1) the \emph{information bottleneck} \cite{tishby1999ib,tishby2015dl}:
$\min_{p(z|s)} I(S;Z)$ s.t.\ $Z$ is value-sufficient; (2) the
\emph{rate--distortion} function $R(D)$ with distortion the value/return
error; (3) the \emph{effective dimension} of the value-relevant manifold
\cite{li2018intrinsic}. VBN's width sweep is a discrete, coarse sampling of
the value rate--distortion curve $R(D)$: the ``monotone'' fingerprint is a
slowly-decaying $R(D)$ (high rate needed), ``flat-high'' is a sharply-decaying
$R(D)$ (low rate suffices). So compressibility is well-defined; VBN is a
practical proxy for it, in the same sense a linear probe
\cite{alain2016probes} is a proxy for representational content.

\paragraph{Is VBN novel?}
Not as an object---and we are explicit about this. Bottlenecking a task head
to measure how much representation a function needs is the shared idea behind
the information bottleneck \cite{tishby1999ib}, probing classifiers
\cite{alain2016probes}, and intrinsic-dimension studies
\cite{li2018intrinsic}. In RL, the \emph{value-equivalence principle}
\cite{grimm2020value} is the direct theoretical cousin: a model need only be
accurate on the information that affects value, and MuZero
\cite{schrittwieser2020muzero} exploits exactly this. State-abstraction theory
\cite{li2006abstraction} and bisimulation metrics \cite{zhang2021dbc} formalize
value-preserving compression of state. Our novelty is not the bottleneck; it
is (i) using value-relevant compressibility as a \emph{predictor of when a
learned WM helps}, and (ii) validating that prediction \emph{cross-model} on
matched tasks. VBN measures the \emph{size} of the value-equivalent subspace;
the paper's claim is that this size governs WM-dependence.

\paragraph{Toward a training-free criterion (answering ``can pure analysis
decide it?'').}
VBN still requires a short probe on a trained checkpoint. A genuinely
\emph{analytical} criterion---predicting WM-dependence from the task
specification alone---appears reachable and is our main open direction. Three
candidate proxies, none requiring the full agent, follow from the same logic:
\begin{itemize}\itemsep2pt
\item \textbf{Reward-propagation horizon.} The WM helps when credit must
travel far. The spectral radius / mixing time of the discounted transition
operator applied to the reward field bounds how many steps of lookahead the
value needs; long horizons (sparse, underactuated \textsc{acrobot}) predict
high WM-dependence, short horizons (dense-reward \textsc{walker}) predict
redundancy---consistent with our data.
\item \textbf{Controllability / observability structure.} From control theory,
underactuated swing-up tasks have ill-conditioned controllability Gramians:
reaching the goal requires composing dynamics the policy cannot shortcut,
which is exactly where a forward model pays off.
\item \textbf{Value rate--distortion from the reward geometry.} $R(D)$ for the
value can in principle be bounded from the smoothness/Lipschitz structure of
$r$ and the transition kernel, giving a closed-form ``flat-high vs.\ monotone''
prediction without training.
\end{itemize}
Whether any of these correlates with the measured WM-dependence gradient across
a larger task suite is a clean, falsifiable next experiment. If it does, the
diagnostic becomes an \emph{a priori} test---decide whether to pay for a world
model before training one.

\section{Related Work}
Model-based RL with learned latent dynamics: Dreamer
\cite{hafner2023dreamerv3}, TD-MPC2 \cite{hansen2024tdmpc2}, MuZero
\cite{schrittwieser2020muzero}. The value-equivalence principle
\cite{grimm2020value} and value-aware model learning motivate models that
preserve only value-relevant information; we supply an empirical instrument
(VBN) for the \emph{size} of that information and connect it to when the model
helps. State abstraction \cite{li2006abstraction} and bisimulation
\cite{zhang2021dbc} formalize value-preserving compression. The information
bottleneck \cite{tishby1999ib,tishby2015dl}, probing \cite{alain2016probes},
and intrinsic dimension \cite{li2018intrinsic} are the representational-analysis
lineage VBN belongs to. Our tasks are standard DeepMind Control
\cite{tassa2018dmc}.

\section{Limitations and Honest Negatives}
We report, rather than hide, three tensions. \textbf{(1)} The prescriptive
gate fails ($n{=}5$); the contribution is diagnostic only. \textbf{(2)} On the
high-dependence anchor (\textsc{acrobot}) the stripped ablation is
high-variance across seeds, so WM-dependence \emph{magnitude} is seed-sensitive
even though its \emph{sign} is not. \textbf{(3)} VBN requires a short training
probe; the training-free criterion of Sec.~\ref{sec:theory} is conjectural. We
also note that \textsc{hopper-vanilla} in Dreamer did not complete on our
hardware (repeated initialization OOM), so the cross-model table uses three
tasks, not four; the fourth (\textsc{hopper-hop}) contributes only its VBN
fingerprint.

\section{Conclusion}
Whether a learned world model helps is a property of the task's
value-information structure, not of the world-model architecture. A cheap
checkpoint-time probe (VBN) measuring value-relevant compressibility predicts
the same WM-dependence ordering across two independent world-model families
and three tasks: essential on \textsc{acrobot}, helpful on \textsc{cheetah},
redundant on \textsc{walker}. The natural prescriptive fix (a plan-time gate)
does not survive replication, sharpening the contribution to a diagnostic. We
ground that diagnostic in the value-equivalence principle and rate--distortion,
and sketch a training-free analytical criterion---reward-propagation horizon,
controllability, and value rate--distortion---that could let practitioners
decide whether a world model has room to help before training one.

\begin{thebibliography}{99}\small
\bibitem{hafner2023dreamerv3} D.~Hafner, J.~Pasukonis, J.~Ba, T.~Lillicrap.
Mastering Diverse Domains through World Models (DreamerV3). \emph{arXiv:2301.04104}, 2023.
\bibitem{hansen2024tdmpc2} N.~Hansen, H.~Su, X.~Wang. TD-MPC2: Scalable,
Robust World Models for Continuous Control. \emph{ICLR}, 2024.
\bibitem{schrittwieser2020muzero} J.~Schrittwieser et al. Mastering Atari, Go,
Chess and Shogi by Planning with a Learned Model (MuZero). \emph{Nature}, 2020.
\bibitem{grimm2020value} C.~Grimm, A.~Barreto, S.~Singh, D.~Silver. The Value
Equivalence Principle for Model-Based Reinforcement Learning. \emph{NeurIPS}, 2020.
\bibitem{tishby1999ib} N.~Tishby, F.~Pereira, W.~Bialek. The Information
Bottleneck Method. \emph{Allerton}, 1999.
\bibitem{tishby2015dl} N.~Tishby, N.~Zaslavsky. Deep Learning and the
Information Bottleneck Principle. \emph{IEEE ITW}, 2015.
\bibitem{alain2016probes} G.~Alain, Y.~Bengio. Understanding Intermediate
Layers Using Linear Classifier Probes. \emph{arXiv:1610.01644}, 2016.
\bibitem{li2018intrinsic} C.~Li, H.~Farkhoor, R.~Liu, J.~Yosinski. Measuring
the Intrinsic Dimension of Objective Landscapes. \emph{ICLR}, 2018.
\bibitem{li2006abstraction} L.~Li, T.~Walsh, M.~Littman. Towards a Unified
Theory of State Abstraction for MDPs. \emph{ISAIM}, 2006.
\bibitem{zhang2021dbc} A.~Zhang, R.~McAllister, R.~Calandra, Y.~Gal, S.~Levine.
Learning Invariant Representations for RL without Reconstruction (DBC).
\emph{ICLR}, 2021.
\bibitem{buckman2018steve} J.~Buckman, D.~Hafner, G.~Tucker, E.~Brevdo, H.~Lee.
Sample-Efficient RL with Stochastic Ensemble Value Expansion (STEVE). \emph{NeurIPS}, 2018.
\bibitem{tassa2018dmc} Y.~Tassa et al. DeepMind Control Suite.
\emph{arXiv:1801.00690}, 2018.
\end{thebibliography}

\end{document}
